\section{The odd-order obstruction}
\label{sec:odd}

\begin{theorem}[Odd obstruction]
\label{thm:odd}
For every odd $k=2m+1\geq3$, every reverse-complement-invariant decycling
set $M$ in $B(q,k)$ satisfies
\[
  |M|\geq N_q(k)+q.
\]
In particular, ordinary minimum cardinality is impossible.
\end{theorem}

\begin{proof}
Let $M$ be such a set. Fix a complementary pair $a,\bar a$ and define the
length-$2m$ lower states
\[
  f_m=(a\bar a)^m,\qquad g_m=(\bar a a)^m.
\]
Both are fixed by reverse complement.  The length-$k$ word
\[
  w_m=(a\bar a)^m\bar a
\]
contains $m$ copies of $a$ and $m+1$ copies of $\bar a$.  If it were a
proper $r$-fold power, then $r>1$ would divide both counts, contradicting
$\gcd(m,m+1)=1$.  Its PCR therefore has full length $k$.

In the lower-state graph this PCR contains the direct arc $g_m\to f_m$ and,
through its other $k-1$ arcs, a directed path from $f_m$ to $g_m$.  Denote
this PCR by $C_a$. Every decycling set meets $C_a$. If $M$ selected exactly
one arc of $C_a$, deleting the direct arc would leave the long path, while
deleting any other arc would leave the direct path. Thus one
endpoint-to-endpoint direction would survive.

By reverse-complement invariance of $M$, that surviving path would have an
image directed from the other fixed endpoint back to the first.
Their concatenation is a nonempty closed directed walk and contains a
directed cycle, a contradiction.

Consequently $|M\cap C_a|\geq2$. The mate PCR $\RC(C_a)$ contains the
same number of selected words. These two PCRs are distinct: their counts of
$a$ and $\bar a$ are interchanged. Repeating the argument for each of the
$q/2$ complementary alphabet pairs gives $q$ distinct PCRs, each requiring
at least two selected words. PCRs from different alphabet pairs are disjoint
because their words use disjoint sets of letters. Every other PCR still
requires at least one selected word. Since PCRs partition $\A^k$,
\[
  |M|\geq 2q+\bigl(N_q(k)-q\bigr)=N_q(k)+q.
\]
\end{proof}

The proof does not determine the exact odd excess; no formula inferred from
small odd orders is used below.
